\subsection{Construction and transfer of the integral calculus}
\label{sec:integral-calculus}
We prove Theorem~\ref{thm:integral-input}. Local special integrals and the general graph are those of Definition~\ref{def:general-graph}.

\begin{lemma}[Measure invariance without a decomposition]
For strongly progressively measurable $X_n,Y$, convergence of the elementary
test errors uniformly over all unit-bounded tests on every finite horizon is
invariant under equivalent probability measures. Explicitly, for every finite
$T$ and $\varepsilon>0$, one $N$ must satisfy
\[
 E_\mu\left[1\wedge\sup_{t\leq T}|(J\cdot X_n)_t-(J\cdot Y)_t|\right]
 \leq\varepsilon
\]
for every $n\geq N$ and every unit-bounded elementary predictable test $J$. Without right continuity,
use the countable horizon skeleton for the supremum.
\end{lemma}
\begin{proof}
Every capped test envelope is measurable and lies in $[0,1]$. If uniform
convergence failed under the target measure, choose $n_k\geq k$ and tests
$J_k$ whose target expectations stay above one positive constant.
Test-uniform convergence under the source measure makes the expectations of
this chosen envelope sequence tend to zero. Pass to convergence in probability,
transfer it by absolute continuity, then use boundedness to recover convergence
of expectations under the target measure, a contradiction. Reverse the
measures for the converse. No special decomposition is involved.
\end{proof}
Associativity of finite sums turns uniform elementary approximation estimates
into this process convergence. Tests only see increments, so initial values
must be fixed separately. Convergence alone does not produce an integrand.

\begin{lemma}[Uniqueness with fixed initial value]
If a strongly progressively measurable sequence converges in this sense to
right-continuous $Y,Z$ with $Y_0=Z_0$ almost surely, then $Y,Z$ are indistinguishable.
\end{lemma}
\begin{proof}
For a fixed test, convergence of capped expectations gives convergence in
probability uniformly on compact intervals, whose limit is unique.
Use the unit test on $(0,T]$ and evaluate at $T$ to obtain
$Y_T-Y_0=Z_T-Z_0$ almost surely. Equality of initial values and right continuity
extend equality from a countable right-dense time set to all times.
\end{proof}

\begin{proposition}[Uniqueness and reduction of measure transfer]
For a fixed source and coefficient the graph gain is unique up to
indistinguishability. Transfer of bounded-coefficient special graphs between
equivalent measures implies transfer of the truncation graph.
\end{proposition}
\begin{proof}
Two certificate sequences have identical coefficients $H^{(n)}$ row by row.
Fixed-source uniqueness identifies their gains. Apply the preceding uniqueness
lemma to this common sequence and its zero-start regular limits.
For measure transfer apply bounded graph transfer separately to each row,
then transfer test-uniform convergence and replace rows by their
indistinguishable representatives. The regular limit and designated gain are
preserved; no special decomposition of the limit is transferred.
\end{proof}
For $|H|\leq1$, every truncation equals $H$. A constant certificate sequence
therefore gives graph membership once zero start and left limits are verified.

\begin{proposition}[Unbounded coefficients and coordinate tails]
Let $H$ have a global integral representation. Every $H^{(n)}$ has a local
representation over the same source. If $h_k$ is the coefficient on completed
coordinate $k$, with controls $q_k,v_k$, then for
$r_{k,n}=h_k1_{\{|h_k|>n+1\}}$,
\[
 \|r_{k,n}\|_{L^1(v_k)}\to0,\qquad \|r_{k,n}\|_{L^2(q_k)}\to0.
\]
The terminal $L^2$ norm of the corresponding completed martingale tail
integral also tends to zero.
\end{proposition}
\begin{proof}
Restrict all coordinates to the same predictable set and glue the compatible
completed coordinates. Each $h_k$ already belongs to both control spaces,
so the large-value tails of its integrable powers vanish. The controls need
not have finite total mass. The terminal martingale norm is the energy norm
restricted to the horizon, bounded by the full coefficient tail norm.
\end{proof}
\begin{proposition}[Uniform truncation convergence and inclusion]
If the source martingale component has left limits, every global or local
special integral representation, including unbounded coefficients, belongs
to the all-time truncation graph with the same coefficient and gain.
\end{proposition}
\begin{proof}
Fix a completed coordinate. For a zero-start square-integrable martingale $M$
and any elementary predictable test with represented integrand $|J|\leq1$,
discrete $L^2$ contraction, chronological grid limits, and Doob's estimate give
\[
 E\left[1\wedge\sup_{t\leq T}|(J\cdot M)_t|\right]\leq2\|M_T\|_2.
\]
This bound does not depend on the number of test intervals or the sum of
coefficient magnitudes. A completed martingale integral is constant after its
horizon, so the estimate applies at any test horizon.
For a variation-integrable coefficient $r$, identify the completed
finite-variation integral $B$ with its pathwise Stieltjes integral. The signed
measure density formula yields
\[
 |dB|(\mathbb R_+)\leq\int|r_t|\,d|A|_t,\qquad
 \sup_t|(J\cdot B)_t|\leq|dB|(\mathbb R_+).
\]
The canonical variation norm bounds the $L^1$ norm of the first right-hand
side. Transfer convergence in probability from the reference measure to the
original probability and use the cap by one to obtain convergence of
expectations, uniformly in the test.
Apply both estimates to complementary coefficient tails. Additivity
identifies the sum of the bounded cut and the tail with the original gain,
so each coordinate has test-uniform truncation convergence.
Choose a localization coordinate first so that $\Pp(\tau_k\leq T)$ is small,
then the row cutoff on that coordinate. Before the stop all tests agree;
after it the capped error is at most one. This removes localization.

For a local representation, stopping at each positive integer horizon gives
a global one. Integrate all bounded truncations on one source schedule and
identify their stops with the global certificates by uniqueness. Remove
localization again. Transfer regular representatives on a countable common
full-measure set and correct its complement using the usual conditions.
This yields an everywhere regular zero-start representative. No common
schedule for the unbounded local coefficient itself is needed.
\end{proof}

\begin{proposition}[Bounded global certificates]
With left limits for the source martingale, a global integral representation
whose coefficient satisfies $|H|\leq b$ belongs to the truncation graph.
The same holds for any bounded-coefficient local representation stopped at a
positive finite deterministic time $T$.
\end{proposition}
\begin{proof}
Restrict each completed coefficient to $\{|H|\leq n+1\}$ and glue the
compatible coordinates. Completed stopping identities realize every finite
stop on the same schedule. For $n+1\geq b$ the coefficients coincide with
$H$, hence their gains are indistinguishable by uniqueness.
Refine the source schedule against itself to obtain martingale coordinates
with left limits. Glue the c\`adl\`ag martingale and locally finite-variation
coordinates into a regular representative of the original gain; the completed
graph gives zero initial value. All test errors on the eventual constant
sequence vanish. For a local representation stop at $T$: its coefficient
$1_{(0,T]}H$ is bounded by $\max(b,0)$ and its gain is $X_{t\wedge T}$.
Apply the global result. Global finite variation of the unstopped gain is
not required.
\end{proof}

\begin{proposition}[Constructing bounded predictable integrals]
The bounded source and any bounded predictable $H$ supply a local integral
graph and a truncation graph. Any existing bounded-coefficient local graph
enters the all-time truncation graph with its designated gain.
\end{proposition}
\begin{proof}
Use the completed schedule constructed directly from the source. Each
coordinate has finite canonical variation and quadratic energy measures.
Boundedness gives $L^1$ and $L^2$ membership, respectively. Use the same $H$
in all coordinates and glue. Refine for left limits, glue the c\`adl\`ag
martingale and locally finite-variation parts, and obtain zero start from the
completed graph. Apply the same construction to every $H^{(n)}$.
For $n+1\geq b$ fixed-source uniqueness makes all gains indistinguishable,
so test errors vanish on one common tail. Uniqueness also identifies any
previously designated local gain with this constructed gain.
\end{proof}

\begin{proposition}[Integral transfer under equivalent measures]
Suppose the same bounded price $S$ has source data under equivalent
probabilities $\mu,\nu$. Bounded predictable local integral graphs transfer
with coefficient and gain unchanged. Consequently the general truncation
graph transfers, and stopped Mémín realizations return to the original measure.
\end{proposition}
\begin{proof}
First take global bounded-coefficient certificates under both measures.
The sums of the two coordinate energy controls and of the two variation
controls are finite measures with no mass at time zero. Common elementary
density gives one sequence $J_n$ approximating the coefficient in both
$L^2$ energy and $L^1$ variation controls. The completed integral errors tend
to zero uniformly over all tests. On each coordinate the elementary integral
is the pathwise $J_n\cdot S$ stopped at that coordinate's stopping time.

Stop further at the minimum of the two coordinate stops. Finite stopping
cannot increase a capped test envelope. Both integrals are thus limits of
one stopped elementary sequence. Transfer convergence between measures and
use zero-start uniqueness to identify these stopped gains. The sequence of
minimum stops is exhaustive, so the stops can be removed. Apply this at each
deterministic horizon for local certificates.
Construct the bounded graph directly under the target source and identify
its gain by this uniqueness argument. Transfer a general graph row by row
through bounded truncations and then transfer convergence to its regular
limit. For a stopped Mémín realization, absolute continuity preserves
exhaustiveness of the stops. No unbounded gain's special decomposition is
transferred in this argument.
\end{proof}

\begin{proposition}[All-time realization by pasting stopping intervals]
If a zero-start regular $X$ has a general truncation graph at every stage of
a finite-valued exhaustive stopping sequence, then one predictable coefficient
$H$ realizes its all-time graph against the original source.
\end{proposition}
\begin{proof}
Place bounded graphs on the source's directly constructed common schedule.
The completed stopping rule gives graphs of arbitrary stopped gains.
Stopping at a finite time preserves convergence against elementary tests,
so the rule passes
to general truncation graphs. For $\sigma\leq\tau$, subtract the stopped
rows to represent coefficient restriction to $(\sigma,\tau]$ and gain
$X^\tau-X^\sigma$. Restriction commutes with coefficient truncation.

The common null set where monotonicity or divergence of the stops fails is
measurable at time zero. Replace stops there by $n+1$, obtaining an everywhere
increasing sequence $\upsilon_n$. This does not change stopped gains up to
indistinguishability. Restrict the first coefficient to $(0,\upsilon_0]$ and
the subsequent ones to $(\upsilon_n,\upsilon_{n+1}]$. Their disjoint supports
make truncation commute with finite sums. Additivity and convergence give
the graph of each finite sum, whose gain telescopes to $X^{\upsilon_n}$.
The coefficients stabilize pointwise on every $(0,\upsilon_n]$; their limit
$H$ is predictable and agrees there with the corresponding finite sum.
Construct its bounded truncation integrals directly from the source, identify
their stops with the finite-sum rows by uniqueness, and remove the localization
error. This gives $(H,X)$ without assuming coefficients agree on overlaps.
\end{proof}

\paragraph{Unit-bounded transforms and Hahn components}
We prove items 4 and 5 of Theorem~\ref{thm:integral-input}. The localization coordinates below have true $L^2$ components on finite horizons.

\begin{proposition}[Completed bounded transforms and finite-horizon Hahn gains]
Let $Y=M+A$ be strongly adapted and right-continuous, with predictable global
bounded-variation $A$ and true martingale $M$ satisfying $M_T\in L^2$.
Suppose every unit elementary test obeys
\[
 E\left[1\wedge\sup_{t\leq T}|(J\cdot M)_t|\right]\leq b.
\]
For predictable $|k|\leq1$, its completed integral
$I^k=(k1_{(0,T]})\cdot M$ has the same test bound $b$.
Moreover a positive predictable Hahn set $B$ gives
$R=1_{B\cap(0,T]}\cdot Y=N+C$ with, almost surely for all $0\leq a\leq c\leq T$,
\[
 0\leq C_c-C_a,\qquad A_c-A_a\leq C_c-C_a.
\]
The component $N$ has the same test bound and
$E[1\wedge\sup_{t\leq T}|N_t|]\leq b$.
\end{proposition}
\begin{proof}
The normalized variation and martingale energy controls are finite.
Predictable ring density gives one elementary sequence with $|J_n|\leq1$
converging to $k1_{(0,T]}$ in both controls. Energy continuity identifies its
martingale integrals and gives elementary-Émery convergence to $I^k$.
For another unit test $J$, finite-horizon associativity gives
$J\cdot(J_n\cdot M)=(JJ_n)\cdot M$ with $|JJ_n|\leq1$.
Pass the bound $b$ through the triangle inequality and arbitrarily small
approximation error. Variation error is unnecessary for the martingale estimate.
Obtain $B$ from the normalized finite-variation bridge and restrict the
horizon's realized unit integral. The completed/Stieltjes identification
expresses $C_c-C_a$ as the signed mass of $B\cap(a,c]$; positivity of Hahn
mass gives both increment bounds, including the right endpoint.
The transform estimate applies to $N$. Its zero start and the horizon unit
test give the maximal-function bound.
\end{proof}

For the local version remove stops as follows.

Use the integral graph's localizers $\sigma_r$ and completion horizons $U_r$.
Its normalized stopped martingale component is the completed integral of
$k$; restriction to $(0,U_r]$ does not change its energy class or unit bound.
Apply the preceding approximation to each coordinate while keeping test
horizon $T$ fixed. Stopping at $U_r$ before testing cannot increase the capped
maximum relative to the test of $M^{\sigma_r}$. The stopped elementary gain
identity bounds the latter by a test of $M$ at the same $T$. No ordering of
$U_r$ and $T$ is required.
On $\{\sigma_r>T\}$ stopped and unstopped tests agree; otherwise the cap
costs at most one. The full expectation is at most
$b+P(\sigma_r\leq T)$; let $r\to\infty$. Centering does not change increments.
A local finite-variation graph may first be stopped at $T+1$.


Elementary density for a general graph follows by approximating each bounded truncation row and diagonalizing over integer horizons. Proposition~\ref{prop:gi-limit} makes its gain a good integrator.
